\documentclass[a4paper,10pt]{article}
\usepackage{hyperref}
\usepackage{paper-en}

\def\thetitle{Tori without contractible closed geodesics}
\def\theauthors{Anton Petrunin}

\hypersetup{colorlinks=true,
citecolor=black,
linkcolor=black,
anchorcolor=black,
filecolor=black,
menucolor=black,
urlcolor=black,
pdftitle={\thetitle},
pdfauthor={\theauthors}}

\begin{document}
\title{\thetitle}
\author{\theauthors}
\date{}
\maketitle

\begin{abstract}
We show that a Riemannian two-torus has no contractible closed geodesics if and only if it admits a geodesic foliation.
\end{abstract}

\section{Introduction}

\begin{thm}{Theorem}\label{thm:main}
A Riemannian 2-torus $(\TT^2,g)$ has no contractible closed geodesics if and only if it admits a geodesic foliation.
\end{thm}

Note that Riemannian 2-torus $(\TT^2,g)$ has a contractible closed geodesic if and only if the univeral metric covering $(\RR^2,\tilde g)$ has a closed geodesic.
Thus the theorem is equivalent to the following statement that we will actually prove.

\begin{thm}{Reformulation}
Let $\tilde g$ is a $\ZZ^2$-periodic Riemannian metric on $\RR^2$.
Then $(\RR^2, \tilde g)$ has no closed geodesics if and only if it admits a $\ZZ^2$-peridic geodesic foliation.
\end{thm}

The if part follows by the argument of Y. Colin de Verdière \cite[7.9]{gromov}, it works for any geodesic foliation, not necessary $\ZZ^2$-peridic.

\parit{Proof of the if part.}
By Poincaré index theorem, a foliation of $\RR^2$ cannot have closed leaves.

Any geodesic in $(\RR^2, \tilde g)$ either runs along the foliation and therefore cannot come back, or it crosses leaves of foliation transversely in one direction and therefore it cannot come back.
\qeds

\section{Preliminaries}\label{sec:prelim}

Let $(\TT^2,g)$ be a Riemannian torus and $(\RR^2,\tilde g)$ its unveral Riemannian covering;
so $\tilde g$ be a $\ZZ^2$-periodic Riemannian metric on $\RR^2$.
We will denote by $|p,q|_{\tilde g}$ the distance induced by $\tilde g$.

\paragraph{Horocompactification.}
Consider the space $C(\RR^2)$ of continuous functions $\RR^2\to \RR$ equipped with the compact-open topology.
Given a point $z\in \RR^2$, consider the function $f_z\colon \RR^2\to \RR$ defined by
\[f_z(x)=|z,x|_{\tilde g}-|z,0|_{\tilde g}.\]
The map $F\colon\RR^2\to C(\RR^2)$ defined by $z\mapsto f_z$ is an embedding with compact closure.
The closure of $F(\RR^2)$ is called \emph{horocompactification} of $(\RR^2, \tilde g)$ \cite{gromov-hyperbolic}.
The points added in the compactification will be called \emph{ideal points} of $(\RR^2, \tilde g)$.

Each ideal point $\xi$ is defined by its function in $C(\RR^2)$, which called Busemann function of $\xi$;
it will be written as $x\mapsto |\xi,x|_{\tilde g}$.
Buseman functions sheres many properties of distance function they can be treated as a normalized distance function to an ideal point.

The set of ideal points is homeomorpic to the circle.

\paragraph{Lines.}
If $f\circ\gamma(t)=\const-t$ for a 1-Lipschitz curve $\gamma$ and 1-Lipschitz function $f$,
then we say that $\gamma$ is \emph{calibrated} by $f$;
in this case $\gamma$ is a minimizing geodesic.


A \emph{line} and \emph{half-line} in a metric space are defined as isometric copy or $\RR$ and $[0,\infty)$, respectively.

Every half-line $h\colon[0,\infty)\to (\RR^2,\tilde g)$ defines a unique Buseman function
\[|\xi,x|_{\tilde g}=\lim_{t\to\infty}(|h(t),x|_{\tilde g}-|h(t),0|_{\tilde g}),\]
which callibrates $h$; in this case we say that $\xi$ is the ideal point of the half-line~$h$.

Converally, for every ideal point $\xi$ and a point $p$ there is a half-line $h$ that starts from $p$ and calibrated by the Buseman function $x\mapsto|\xi,x|_{\tilde g}$.
Such a half-line may not be unique.
Thus the ideal points can be defined as equivalence class of half-lines in $(\RR^2,\tilde g)$ calibrated by single Buseman function.

\paragraph{Ideal directions}
By Hedlund's theorem \cite{hedlund,burago} half-lines in $(\RR^2,\tilde g)$ lie on a bounded distance from a Euclidean half-line (it also holds for Finsler metrics \cite[3.1]{schroeder-survey}).
In particular, we can define asymptotic direction $\dir h\in\SSS^1$ of any half-line $h$ in $(\RR^2,\tilde g)$.
For a line $\ell$ we have two asymptotic direction of entrance and exit $\dir^-\ell$ and $\dir^+\ell$, which are opposite by Hedlund's theorem, so $\dir^-\ell=-\dir^+\ell$.

Given a direction $u\in \SSS^1$ and a point $p\in\RR^2$ we can find a half-line $h\colon[0,\infty)\to(\RR^2,\tilde g)$ with $\dir h=u$ with an ideal point $\xi$.
If $u$ is irratonal, then $\xi$ is uniquely defined; we will write $\xi=\bar u$.
In this case any $\ZZ^2$-shift of $x$ changes $|\xi,x|$ by a constant \cite[8.7(i)]{schroeder-survey}.

A rational direction $u\in \SSS^1$ may correspond to a nontrivial closed segment of ideal point, let us denote by
$\cev u$ and $\vec{u}$ the left and right ends.
If $\cev u=\vec u$, then $\TT^2$ is folliated by closed geodesic with asymptotic direction $u$.

If line $\ell$ has irrational direction $u$, then it is calibrated by $x\mapsto |\bar u,x|$ and $x\mapsto -|-\bar u,x|$.
The set of all such lines are disjoint;
their union is the set of minimum of the sum
\[x\mapsto |\bar u,x|+|{-}\bar u,x|.\]
This function is $\ZZ^2$-periodic;
indeed, $\ZZ^2$-shift changes the sum by a constant and by the triangle inequality, the functions is bounded below.
If $u$ is rational, then the same works for the sums
\[x\mapsto |\cev u,x|+|{-}\cev u,x|
\quad\text{and}\quad
x\mapsto |\vec u,x|+|{-}\vec u,x|.\]

If $u_n$ is a sequence of irrational directions such that $\bar u_n\to \cev u+$ or $\bar u_n\to \vec u-$,
then $|\bar u_n,x|\to|\cev u,x|$ or $|\bar u_n,x|\to|\vec u,x|$, respectively \cite[4.15]{schroeder-survey}.



\paragraph{Strips.}
The following statements are covered in \cite[§ 4--5]{bangert-gutkin}.

A line in $(\RR^2,\tilde g)$ will be called \emph{periodic} if it is invariant with respect to some $\ZZ^2$ sift;
these lines project to periodic geodesics in the torus.
If $\ell$ is periodic line for a shift in direction $u$ (which must be rational) then $\ell$ is calibrated by all four Busemann functions $x\mapsto \pm|\pm\cev u,x|$ and $x\mapsto \pm|\pm\vec u,x|$.

Applying an $\SL(2,\ZZ)$-matrix to the plane, we may assume that given the $\ZZ^2$-sift is $(x,y)\mapsto (x+1,y)$.
In this case, periodic lines are called \emph{horizontal}.
These are disjont lines.
Their union is a closed set, say $H$ and the complement $\RR^2\setminus H$ is a union of disjoint introiors of   strips between pairs of horizontal lines.

Let $\Sigma$ be a strip between lower and upper horizontal lines $\ell$ and $m$, and let $\xi=(1,0)$.
Denote their ideal points by $\lambda^\pm$ and $\mu^\pm$.
The restrictions of their Buseman functions can be defined within $\Sigma$;
up to a constant, these restrictions coinside with $x\mapsto \pm|\pm\cev u,x|$ and $x\mapsto \pm|\pm\vec u,x|$.

The Buseman functions $x\mapsto |\cev u,x|$ and $x\mapsto -|{-}\cev u,x|$ calibrate that come in North-East direction in $\Sigma$ they have left asymptate at $\ell$ and right asymptotes on $m$.
Similarly, $x\mapsto |\vec u,x|$ and $x\mapsto -|{-}\vec u,x|$ calibrate that come in Sout-East direction in $\Sigma$ they have left asymptate at $m$ and right asymptotes on $\ell$.


%???PIC



\section{Full turn from small turns}

Let $s$ is a geodesic in $(\RR^2,\tilde g)$ with minimizing tales;
that is, both tales of $s$ are half-lines, so the asymptotic direction of entrance and exit $\dir^-s$ and $\dir^+s$ are defined.
We say that $s$ \emph{turns across} $\xi\in\SSS^1$ if $\xi$ lies on the left from $-\dir^-s$ and on the right $\dir^+s$;
more precisely, the directions $\dir^-s$, $\dir^+s$, $\xi$ and $-\dir^-s$ appear in the same counterclockwise  cyclic order on $\SSS^1$.

\begin{thm}{Key lemma}\label{lem:key}
Suppose that the plane $\RR^2$ equipped with a $\ZZ^2$-periodic Riemannian metric $\tilde g$ does not admit a line foliation.
Then for any direction $\xi$ of $\RR^2$ there is a geodesic $s$ with minimizing tales that turns across $\xi$.

\end{thm}

\parbf{Proof of the main theorem modulo the key lemma.}
We can choose finite number of geodesic $s_1,\dots,s_n$ that turn across all directions.
Since these geodesics have minimizing tales, the quoted Hedlund's theorem implies that each $s_i$ lies on a bounded distance from a Euclidean angle, say $r_i$ (made from two Euclidean rays).

Construct a polygon $P$ from translated arcs of $r_i$ so that all angles at the join of arc are concave; that is, bigger than $\pi$.
So the only convex vertices of $P$ are the corners inside the translated arcs.
Rescaling $P$ up makes it possible to approximate all the arcs by integer translations of arcs of $s_i$.
Since a pair of minimizing tales may cross only once,
this produces a geodesic polygon $Q$ in $(\RR^2,\tilde g)$ with all concave angles, so the closure of the complement of $Q$ is locally convex.

It remains to choose a shortest curve $\gamma$ that surrounds $Q$ and observe that $\gamma$ is a closed geodesic in $(\RR^2,\tilde g)$.
\qeds

\section{Gates}

The following observation provides a tool for constructing geodesic in the key lemma.
Intuitively speaking, we scratch curves that pass a gate or two; if the obtained curve does not touch gate sides, then we get a geodesic.

\begin{thm}{Observation}\label{obs:join}
Let $\tilde g$ be $\ZZ^2$-periodic Riemannian metric on $\RR^2$,
and let $z$ be a periodic line that separates two irrational directions $\alpha,\omega\in\SSS^1$.
There is a turn from $\alpha$ to $\omega$ if one of the following condition holds:

\begin{subthm}{}
There is an arc $I$ in $z$ such that the minimum of function
\[f\colon p\mapsto \bus_\alpha(p)+\bus_\omega(p)\]
for $p\in I$ is strictly below both endpoint values.
\end{subthm}

%PIC???

\begin{subthm}{}\label{obs:main2}
There are two arcs $I$ in $z$ and $J$ on its $\ZZ^2$-shift $z'$ to the side of $z$ with $\omega$
and the function
\[f\colon (p,q)\mapsto \bus_\alpha(p)+|p-q|_g+\bus_\omega(q)\]
for $p\in I$ and $q\in J$
is strictly below its boundary values.
\end{subthm}

%PIC???

\end{thm}

\parbf{Proof of the key lemma.}
Suppose $\xi\in\SSS^1$ is an irrational direction.
Choose a periodic line $z$ that separates $-\xi$ and $\xi$.
Note that $z$ crosses lines that calibrated by $\bus_\xi$.
But not all points on $z$ lie on such lines, so $\ex_\xi(p)>0$ for some points $p\in z$.

Thus we can choose a segment $I$ in $z$ that meets the condition of the observation for the function
$f=\bus_{-\xi}+\bus_\xi$.
Choose $\alpha\approx -\xi$ and $\omega\approx \xi$ and note that $I$ meets the condition for the function
$f=\bus_{\alpha}+\bus_\omega$, and the observations finishes the proof.

Now suppose $\xi$ is a rational direction.
We may assume that $\xi$ corresponds to the horizontal shift $\tau\colon(x,y)\mapsto (x+1,y)$.
Applying the observations about horizontal lines and strips at the end of §~\ref{sec:prelim},
we can a strip $\Sigma$ bounded by horizontal lines $\ell$ and $u$ such that $\ex_{\xi\downarrow}$ and $\ex_{\xi\uparrow}$ are not constant in $\Sigma$.
Note that we can choose a gate $I$ near $\ell$ that meet the first condition for function $f=\bus_{-\xi\uparrow}+\bus_{\xi\uparrow}$.
The same way we can choose a gate $J$ near $\ell$ that meet the first condition for function $f=\bus_{-\xi\downarrow}+\bus_{\xi\downarrow}$.
Now let us apply the huge horizontal shift to move $J$ far to the right from $I$.
Minimizing geodesics between $I$ and $J$ nearly approach $\ell$ in the middle.
And we can conclude that the pair $I$ and $J$ meet the condition in \ref{obs:main2} for the function
\begin{align*}
f\colon (p,q)&\mapsto \bus_{-\xi\uparrow}(p)+|p-q|_{\tilde g}+\bus_{\xi\downarrow}(q).
\intertext{Therefore they also meet the assumptions of the observation for the function}
f\colon (p,q)&\mapsto\bus_{\alpha}(p)+|p-q|_{\tilde g}+\bus_{\omega}(q)
\end{align*}
if
$\alpha$ close to $-\xi$ and lying on its right side
and
$\omega$ close to $\xi$ and lying on its left side.
Again the observation finishes the proof.
\qeds




\parbf{Acknowledgments.}
The proof was found in collaboration with AI (GPT-6.2 Astra).

{\sloppy
\def\emph{\textit}
\printbibliography[heading=bibintoc]\fussy}
\end{document}
